\documentclass[11pt,a4paper]{article}

% ============================================================
% PAQUETES BÁSICOS
% ============================================================
\usepackage{fontspec}
\setmainfont{Myriad Pro}    % Fuente principal
\usepackage[spanish]{babel}
\usepackage{graphicx}
\usepackage{geometry}
\usepackage{fancyhdr}
\usepackage{titlesec}
\usepackage{xcolor}
\usepackage{hyperref}
\usepackage{amsmath}        % Para ecuaciones matemáticas
\usepackage{float}          % Para mejor control de posición de figuras
\usepackage{caption}        % Para personalizar captions
\usepackage{subcaption}     % Para subfiguras si se necesitan
\usepackage{microtype}      % Mejora el espaciado y reduce overfull
\usepackage{eso-pic}        % Paquete clave para añadir imágenes de fondo
% Paquetes recomendados para tablas
\usepackage{placeins}       % \FloatBarrier para controlar floats
\usepackage{booktabs}       % Reglas profesionales (\toprule, \midrule, \bottomrule)
\usepackage{tabularx}       % Columnas X que ajustan al ancho
\usepackage{array}          % Control avanzado de columnas
\newcolumntype{L}{>{\raggedright\arraybackslash}X} % Columna X alineada a la izquierda


% ============================================================
% CONFIGURACIONES BÁSICAS
% ============================================================

% Configuración de márgenes
\geometry{
    top=2.5cm,
    bottom=2.5cm,
    left=2cm,
    right=2cm,
    headheight=14.5pt
}

% Configuración de enlaces
\hypersetup{
    colorlinks=true,
    linkcolor=blue,
    filecolor=magenta,
    urlcolor=cyan,
    pdftitle={Memoria Práctica},
    pdfpagemode=FullScreen,
}

% Configuración de encabezados y pies de página
\pagestyle{fancy}
\fancyhf{}
\fancyhead[L]{Sistemas Operativos - Práctica 3}
\fancyhead[R]{Diciembre - 2025}
\fancyfoot[R]{\thepage}
\fancyfoot[L]{Julian Hinojosa}

% ============================================================
% CONFIGURACIONES ESPECÍFICAS POR ASIGNATURA
% ============================================================

% --- [SISTEMAS OPERATIVOS] Configuración de listings para código ---
\usepackage{listings}

% Configuración de colores para el código
\definecolor{codegreen}{rgb}{0,0.6,0}
\definecolor{codegray}{rgb}{0.5,0.5,0.5}
\definecolor{codepurple}{rgb}{0.58,0,0.82}
\definecolor{backcolour}{rgb}{0.95,0.95,0.92}
\definecolor{codeorange}{rgb}{0.8,0.4,0}
\definecolor{framecolor}{rgb}{0.7,0.7,0.7}

% Configuración de listings para código C
\lstdefinestyle{mystyle}{
    backgroundcolor=\color{backcolour},   
    commentstyle=\color{codegreen}\itshape,
    keywordstyle=\color{blue}\bfseries,
    numberstyle=\tiny\color{codegray},
    stringstyle=\color{codepurple},
    basicstyle=\ttfamily\small,
    breakatwhitespace=false,         
    breaklines=true,                 
    captionpos=b,                    
    keepspaces=true,                 
    numbers=left,                    
    numbersep=8pt,
    showspaces=false,                
    showstringspaces=false,
    showtabs=false,                  
    tabsize=4,
    frame=single,
    frameround=tttt,
    rulecolor=\color{framecolor},
    framesep=4pt,
    xleftmargin=15pt,
    xrightmargin=5pt,
    language=C,
    extendedchars=true,
    inputencoding=utf8,
    escapeinside={(*@}{@*)},
    morecomment=[l][\color{codeorange}]{\#},
    columns=flexible,
    aboveskip=15pt,
    belowskip=10pt,
    literate={á}{{\'a}}1 {é}{{\'e}}1 {í}{{\'i}}1 {ó}{{\'o}}1 {ú}{{\'u}}1
             {Á}{{\'A}}1 {É}{{\'E}}1 {Í}{{\'I}}1 {Ó}{{\'O}}1 {Ú}{{\'U}}1
             {ñ}{{\~n}}1 {Ñ}{{\~N}}1 {¿}{{?`}}1 {¡}{{!`}}1
}

\lstset{style=mystyle}

% Configuración para el caption de los listings
\DeclareCaptionFont{white}{\color{white}}
\DeclareCaptionFormat{listing}{\colorbox{gray}{\parbox{\dimexpr\textwidth-2\fboxsep\relax}{#1#2#3}}}
\captionsetup[lstlisting]{format=listing, labelfont=white, textfont=white, singlelinecheck=false, margin=0pt, font={bf,footnotesize}}
% --- FIN [SISTEMAS OPERATIVOS] ---

\begin{document}

% --- Portada y índice sin numeración ---
\pagenumbering{gobble} % no mostrar números de página en preliminares

% Página de portada - Imagen ocupando la totalidad de la página
\begin{titlepage}
    % Esto pone la imagen detrás de todo, ocupando toda la hoja
    \AddToShipoutPictureBG*{%
        \AtPageLowerLeft{%
            % Asegúrate que la imagen se llama 'portada.png'
            \includegraphics[width=\paperwidth,height=\paperheight]{portada.png}%
        }%
    }
    \null % Comando invisible para que LaTeX sepa que aquí hay una página
\end{titlepage}

\newpage

% Tabla de contenidos
\tableofcontents
\newpage

% --- A partir de aquí numeración normal ---
\pagenumbering{arabic}   % cambiar a 1,2,3...
\setcounter{page}{1}     % empezar en 1
\pagestyle{fancy}        % restaurar encabezados/pies definidos en el preámbulo

% Inicio del documento
\section{Introducción}

Una de las funciones más importantes de los sistemas operativos es gestionar y organizar el sistema de memoria. En esta práctica se implementa un simulador de gestión de memoria basado en el esquema de **particiones dinámicas o variables**.

El programa desarrollado, denominado \texttt{gestormemoria}, permite simular la carga y descarga de procesos en un espacio de memoria de 2000 unidades, utilizando una unidad de asignación mínima de 100. Se han implementado dos algoritmos de asignación de particiones fundamentales:
\begin{itemize}
    \item \textbf{Primer hueco} (First Fit): Asigna el primer bloque libre que sea suficiente para el proceso.
    \item \textbf{Siguiente hueco} (Next Fit): Similar al anterior, pero comienza la búsqueda a partir de la última asignación realizada.
\end{itemize}

La implementación se ha realizado en lenguaje C, incluyendo la generación de un fichero de registro (\texttt{particiones.txt}) y una visualización gráfica del estado de la memoria (utilizando Raylib) para facilitar la comprensión del funcionamiento de los algoritmos.

\section{Objetivos}

Los objetivos de esta práctica son:

\begin{itemize}
    \item Comprender los conceptos fundamentales de la gestión de memoria mediante asignación contigua y particiones dinámicas.
    \item Implementar las estructuras de datos necesarias para gestionar particiones de memoria y huecos libres.
    \item Programar y comparar el funcionamiento de los algoritmos de asignación \textit{primer hueco} y \textit{siguiente hueco}.
    \item Simular la ejecución temporal de procesos, gestionando su llegada, asignación de memoria y liberación de recursos tras su ejecución.
    \item Visualizar gráficamente el estado de la memoria para analizar la fragmentación y el comportamiento de los algoritmos.
\end{itemize}

\section{Desarrollo}

El código fuente del proyecto se ha estructurado de forma modular para separar la lógica de la simulación, la gestión de entrada/salida y la interfaz 
de usuario. A continuación se describen los principales componentes y archivos que conforman la solución:

\subsection{Estructura del Proyecto}

\begin{itemize}
    \item \textbf{main.c}: Es el punto de entrada de la aplicación. Su función principal es inicializar la simulación y crear dos procesos mediante 
    la llamada al sistema \texttt{fork()}:
    \begin{itemize}
        \item \textbf{Proceso Padre}: Ejecuta la interfaz textual (TUI) para depuración y control paso a paso por consola.
        \item \textbf{Proceso Hijo}: Inicializa la librería \textbf{Raylib} y gestiona la interfaz gráfica (GUI) para la visualización en tiempo real
        de la memoria y el estado de los procesos.
    \end{itemize}
    
    \item \textbf{sim\_engine.c / .h}: Constituye el núcleo o motor de la simulación. Define las estructuras de datos fundamentales (\texttt{Memoria}, 
    \texttt{Particion}, \texttt{Proceso}) e implementa la lógica de gestión de memoria:
    \begin{itemize}
        \item Inicialización del sistema de particiones.
        \item Algoritmos de búsqueda de huecos: \textit{Primer Hueco} y \textit{Siguiente Hueco}.
        \item Asignación de memoria (con división de huecos si es necesario) y compactación automática (fusión de huecos adyacentes al liberar un proceso).
        \item Control del ciclo de vida de los procesos (llegada, ejecución/envejecimiento y finalización) a través de la función \texttt{avanzar\_tiempo}.
    \end{itemize}
    
    \item \textbf{ficheros.c / .h}: Módulo encargado de la persistencia y entrada de datos. Para garantizar la robustez en las operaciones de entrada/salida, este módulo se apoya en la librería externa \textbf{file\_utils} (ubicada en la raíz del repositorio, en el directorio \texttt{lib/}).
    \begin{itemize}
        \item \texttt{cargar\_procesos}: Lee y parsea el archivo de entrada con la definición de los procesos utilizando llamadas al sistema POSIX y las utilidades de lectura completa de la librería mencionada.
        \item \texttt{guardar\_estado}: Escribe el estado actual de las particiones en el fichero de log (\texttt{particiones.txt}) en cada instante 
        de la simulación.
    \end{itemize}
    \textit{Nota sobre dependencias}: Debido al uso de esta librería de utilidades compartida (\texttt{lib/file\_utils.h}), es imprescindible disponer de la estructura completa del repositorio para la correcta compilación del proyecto.
\end{itemize}

\subsection{Estructuras de Datos Principales}

Para modelar el problema se han definido las siguientes estructuras en \texttt{sim\_engine.h}:

\begin{itemize}
    \item \textbf{Proceso}: Almacena la información de cada tarea, incluyendo su nombre, tiempos de llegada y ejecución, memoria requerida (que será alineada)
    y banderas de estado (en memoria, finalizado).
    \item \textbf{Particion}: Representa un bloque de memoria contigua. Puede estar en estado ocupado (asignado a un proceso) o libre (hueco). Contiene la 
    dirección de inicio y su tamaño.
    \item \textbf{Memoria}: Estructura contenedora que mantiene el array de particiones, el contador de particiones activas y el puntero al último índice
    asignado para soportar el algoritmo \textit{Siguiente Hueco}.
\end{itemize}

\newpage

\section{Código Fuente Destacado}

A continuación se presentan los fragmentos más relevantes de la implementación, correspondientes a la lógica del motor de simulación (\texttt{sim\_engine.c}).

\subsection{Algoritmos de Búsqueda de Huecos}

Esta función implementa las dos estrategias solicitadas. Para \textit{Primer Hueco} se realiza una búsqueda lineal desde el inicio, mientras que para \textit{Siguiente Hueco} se utiliza una búsqueda circular comenzando desde \texttt{ultimo\_indice\_asignado}.

\begin{lstlisting}[caption={Implementación de First Fit y Next Fit}, label={lst:buscar_hueco}]
int buscar_hueco(Memoria *m, int mem_requerida, TipoAlgo tipo_algo) {
    // Algoritmo Primer Hueco
    if (tipo_algo == ALGO_PRIMER_HUECO) {
        // Recorremos todas las particiones buscando el primer hueco que quepa
        for (int i = 0; i < m->cant_particiones; i++) {
            if (m->particiones[i].estado == 0 && m->particiones[i].tamano >= mem_requerida) {
                return i; // Devolvemos el índice del hueco encontrado
            }
        }
    }

    // Algoritmo Siguiente Hueco
    if (tipo_algo == ALGO_SIGUIENTE_HUECO) {
        int inicio = m->ultimo_indice_asignado;
        int total = m->cant_particiones;

        // Búsqueda circular comenzando desde la última asignación
        for (int j = 0; j < total; j++) {
            int i = (inicio + j) % total;
            if (m->particiones[i].estado == 0 && m->particiones[i].tamano >= mem_requerida)
                return i;
        }
    }
    return -1;
}
\end{lstlisting}

\subsection{Compactación de Memoria}

Para evitar la fragmentación externa excesiva, se implementa una función de compactación que fusiona huecos adyacentes. Esta función es invocada automáticamente tras la liberación de un proceso.

\begin{lstlisting}[caption={Función de Compactación de Huecos}, label={lst:compactar}]
void compactar(Memoria *m) {
    // Recorremos hasta el penúltimo elemento
    for (int i = 0; i < m->cant_particiones - 1; i++) {
        // Si la partición actual y la siguiente son huecos, las unimos
        if (m->particiones[i].estado == 0 && m->particiones[i+1].estado == 0) {
            // 1. Sumamos tamaños de las particiones
            m->particiones[i].tamano += m->particiones[i+1].tamano;
            
            // 2. Movemos todas las particiones siguientes una posición a la izquierda
            for (int j = i + 1; j < m->cant_particiones - 1; j++) {
                m->particiones[j] = m->particiones[j+1];
            }
            
            // 3. Actualizamos el contador y retrocedemos índice para re-verificar
            m->cant_particiones--;
            i--;
        }
    }
}
\end{lstlisting}

\subsection{Motor de Tiempo y Ciclo de Vida}

La función \texttt{avanzar\_tiempo} orquesta la simulación en cada "tick" de reloj, gestionando el envejecimiento de los procesos en ejecución, la liberación de memoria al terminar y la admisión de nuevos procesos.

\begin{lstlisting}[caption={Lógica principal de avance temporal}, label={lst:avanzar}]
void avanzar_tiempo(Memoria *m, Proceso procesos[], int num_procesos, int *reloj_actual, TipoAlgo algo, const char* ruta_log) {
    // 1. Envejecimiento y finalización
    for (int i = 0; i < num_procesos; i++) {
        if (procesos[i].en_memoria && !procesos[i].finalizado) {
            procesos[i].t_restante--;

            if (procesos[i].t_restante <= 0) {
                liberar_proceso(m, procesos[i].nombre);
                procesos[i].en_memoria = false;
                procesos[i].finalizado = true;
            }
        }
    }

    // 2. Llegada de nuevos procesos (Planificador a largo plazo)
    for (int i = 0; i < num_procesos; i++) {
        if (!procesos[i].en_memoria && !procesos[i].finalizado && procesos[i].t_llegada <= *reloj_actual) {
            if (asignar_proceso(m, procesos[i], algo)) {
                procesos[i].en_memoria = true;
                procesos[i].t_restante = procesos[i].t_ejecucion;
            }
        }
    }

    // 3. Guardar estado y avanzar reloj
    guardar_estado(ruta_log, m, *reloj_actual);
    (*reloj_actual)++;
}
\end{lstlisting}

\subsection{Integración con File Utils}

Para asegurar la integridad en las operaciones de entrada/salida frente a posibles interrupciones del sistema o transferencias parciales, se hace uso de las funciones \texttt{read\_all} y \texttt{write\_all} proporcionadas por la librería auxiliar. A continuación se muestra cómo se utilizan para la carga segura del archivo de definición de procesos en \texttt{ficheros.c}:

\begin{lstlisting}[caption={Lectura robusta del fichero de entrada}, label={lst:file_utils_usage}]
int cargar_procesos(const char* ruta, Proceso procesos[]) {
    // 1. Abrir descriptor de archivo (POSIX)
    int fd = open(ruta, O_RDONLY);
    if (fd == -1) {
        perror("Error abriendo fichero");
        return 0;
    }

    char buffer_archivo[SIZE_BUFFER_LECTURA];
    memset(buffer_archivo, 0, SIZE_BUFFER_LECTURA);

    // 2. Uso de read_all para garantizar lectura completa
    // Maneja internamente bucles por lecturas parciales o EINTR
    ssize_t leidos = read_all(fd, buffer_archivo, SIZE_BUFFER_LECTURA - 1);
    
    close(fd); // Cierre seguro tras lectura

    if (leidos == -1) return 0;

    // 3. Procesamiento del buffer (parsing)...
    // ...
    return count;
}
\end{lstlisting}

De manera análoga, la función \texttt{guardar\_estado} utiliza \texttt{write\_all} para asegurar que las líneas de log se escriban completas en el disco, evitando escrituras fragmentadas.

\section{Instalación y Ejecución}

Dado que se entrega el repositorio completo del proyecto, a continuación se detallan los pasos necesarios para compilar y ejecutar el simulador en distintos 
sistemas operativos.

\subsection{Requisitos Previos}

El proyecto requiere un compilador de C, CMake y la librería Raylib.

\subsubsection{macOS}
\begin{itemize}
    \item Instalar herramientas de desarrollo: \texttt{xcode-select --install}
    \item Instalar CMake y Raylib (vía Homebrew):
    \begin{lstlisting}[language=bash]
brew install cmake raylib
    \end{lstlisting}
\end{itemize}

\subsubsection{Linux (Debian/Ubuntu)}
\begin{itemize}
    \item Instalar compilador, CMake y dependencias gráficas:
    \begin{lstlisting}[language=bash]
sudo apt update
sudo apt install build-essential cmake libraylib-dev
    \end{lstlisting}
\end{itemize}

\subsubsection{Windows}
Para la instalación de Raylib en Windows, se recomienda consultar la guía oficial en el siguiente enlace. En general, para facilitar el proceso de instalación y configuración de entornos de desarrollo como este, se sugiere el uso de sistemas operativos UNIX-like (macOS o Linux).

\begin{itemize}
    \item Guía detallada para trabajar con Raylib en Windows: \url{https://github.com/raysan5/raylib/wiki/Working-on-Windows}
\end{itemize}

\subsection{Compilación}

El proceso de construcción es estándar gracias a CMake. Asegúrese de estar en el directorio `practica3` (ubicado en la raíz del repositorio en 
`noob-code/UA/year-2/OS/practica3`) y ejecute los siguientes comandos en su terminal:

\begin{lstlisting}[language=bash, caption={Comandos de construcción}]
mkdir build
cd build
cmake ..
make
\end{lstlisting}

\textit{Nota para Windows (MinGW)}: Si el comando \texttt{make} no está disponible, utilice \texttt{mingw32-make}.

Esto generará el ejecutable \texttt{gestomemoria} (o \texttt{gestomemoria.exe} en Windows) dentro del directorio \texttt{build}.

\subsection{Ejecución}

Asegúrese de que el archivo de datos \texttt{entrada.txt} se encuentra en el mismo directorio que el ejecutable o que el programa pueda acceder a él.

\begin{lstlisting}[language=bash, caption={Ejecución del programa}]
./gestomemoria
\end{lstlisting}

Al iniciar, el programa lanzará automáticamente:
\begin{enumerate}
    \item Una ventana gráfica (Raylib) para visualizar la memoria.
    \item La terminal original para ver los logs detallados y mensajes de depuración.
\end{enumerate}

\section{Conclusiones}

La realización de esta práctica ha permitido consolidar la comprensión de los mecanismos fundamentales de la gestión de memoria en sistemas operativos, 
con un enfoque particular en las particiones dinámicas. Hemos podido observar cómo la implementación de algoritmos como Primer Hueco y Siguiente Hueco 
influye directamente en la eficiencia de la asignación de memoria y en la aparición de fragmentación.

Se ha demostrado la importancia de la compactación de memoria para mitigar la fragmentación externa, un problema inherente a la gestión de particiones 
variables. La simulación visual mediante Raylib ha sido una herramienta invaluable para entender de forma intuitiva el comportamiento dinámico de la memoria, 
la creación y fusión de huecos, y el ciclo de vida de los procesos.

Finalmente, la experiencia de desarrollar un simulador desde cero, gestionando aspectos como la carga de procesos, la asignación, liberación y el avance 
temporal, ha proporcionado una perspectiva práctica y profunda sobre los desafíos y soluciones en la administración de recursos en un sistema operativo. 
Se ha logrado un sistema robusto que permite la experimentación y análisis de diferentes escenarios de carga y algoritmos de gestión.

\section{Bibliografía}

\begin{itemize}
  \item \textbf{Transparencias de la Asignatura}: Tema 4 - Gestión de Memoria. Sistemas Operativos (Grado en Ingeniería Informática), Universidad de Alicante, 2025-2026.
  \item \textbf{Enunciado de la Práctica}: Práctica 3 - Gestión de memoria. Sistemas Operativos, Universidad de Alicante, 2025/2026.
  \item \textbf{Documentación de C}: \textit{cppreference.com - C reference}. Disponible en: \url{https://en.cppreference.com/w/c}.
  \item \textbf{Documentación de Raylib}: \textit{Raylib Cheatsheet \& Wiki}. Disponible en: \url{https://www.raylib.com/cheatsheet/cheatsheet.html}.
  \item \textbf{POSIX Standard}: \textit{The Open Group Base Specifications Issue 7}. (Para referencias sobre \texttt{fork}, \texttt{open}, \texttt{read}).
\end{itemize}
  
\end{document}